% --- Archon physics formalization source begin ---
% archon:physics
% archon:covers IPhO2026Problems/problem_IPhO_2026_1_A_1.lean
% archon:source-report reports/ipho_2026/problem_IPhO_2026_1_A_1.source.json
% archon:problem-id IPhO_2026_1
% archon:part-id A.1

\chapter{Physics problem IPhO\_2026\_1\_A\_1}
\label{ch:IPhO2026Problems_problem_IPhO_2026_1_A_1}

\paragraph{Problem source.}
Two water reservoirs are separated by a vertical wall MN.  A square slot of
vertical size a*sqrt(2)/2 is sealed by a fully submerged solid cube of side a
and density 3*rho\_0, where rho\_0 is the density of water.  The cube is hinged
frictionlessly at O and may rotate about an axis perpendicular to the figure.
The maximum permitted difference in water levels is Delta h = 1.41 m.  Use
Figure 1a on the source page for the exact geometry and lever arms.

Current subquestion:
Calculate the side length a that makes Delta h = 1.41 m the maximum permissible water-level difference.

\paragraph{Current subquestion.}
Calculate the side length a that makes Delta h = 1.41 m the maximum permissible water-level difference.

\paragraph{Recorded answer/context.}
a = Delta h/(2*sqrt(2)) = 0.50 m.

\paragraph{Figure/image path.}
/root/proposal\_for\_physic/science-mango/ipho\_2026\_source/image/T1\_page-1.png

\paragraph{Formalization target.}
create a compiling Lean file with sorry bodies at `IPhO2026Problems/problem\_IPhO\_2026\_1\_A\_1.lean`. The Lean declarations must preserve the physical quantities, dimensions or dimensional roles, figure labels, governing-law hypotheses, and final relation expressed by this problem.
Use Mathlib/Physlib names found through LeanExplore where available. If a physics API is missing, introduce faithful local abstractions rather than scalar placeholder aliases.

\begin{definition}[LengthSI]
  \label{decl:physics:IPhO_2026_1_A_1:LengthSI}
  \lean{IPhO2026Problems.IPhO2026_1_A_1.LengthSI}
  A length read in metres.
\end{definition}
\begin{proof}
  This is the defining declaration; unfolding it gives the stated typed data or relation.
\end{proof}

\begin{definition}[AreaSI]
  \label{decl:physics:IPhO_2026_1_A_1:AreaSI}
  \lean{IPhO2026Problems.IPhO2026_1_A_1.AreaSI}
  An area read in square metres.
\end{definition}
\begin{proof}
  This is the defining declaration; unfolding it gives the stated typed data or relation.
\end{proof}

\begin{definition}[VolumeSI]
  \label{decl:physics:IPhO_2026_1_A_1:VolumeSI}
  \lean{IPhO2026Problems.IPhO2026_1_A_1.VolumeSI}
  A volume read in cubic metres.
\end{definition}
\begin{proof}
  This is the defining declaration; unfolding it gives the stated typed data or relation.
\end{proof}

\begin{definition}[MassDensitySI]
  \label{decl:physics:IPhO_2026_1_A_1:MassDensitySI}
  \lean{IPhO2026Problems.IPhO2026_1_A_1.MassDensitySI}
  A mass density read in kilograms per cubic metre.
\end{definition}
\begin{proof}
  This is the defining declaration; unfolding it gives the stated typed data or relation.
\end{proof}

\begin{definition}[AccelerationSI]
  \label{decl:physics:IPhO_2026_1_A_1:AccelerationSI}
  \lean{IPhO2026Problems.IPhO2026_1_A_1.AccelerationSI}
  An acceleration read in metres per second squared.
\end{definition}
\begin{proof}
  This is the defining declaration; unfolding it gives the stated typed data or relation.
\end{proof}

\begin{definition}[PressureSI]
  \label{decl:physics:IPhO_2026_1_A_1:PressureSI}
  \lean{IPhO2026Problems.IPhO2026_1_A_1.PressureSI}
  A pressure read in pascals.
\end{definition}
\begin{proof}
  This is the defining declaration; unfolding it gives the stated typed data or relation.
\end{proof}

\begin{definition}[ForceSI]
  \label{decl:physics:IPhO_2026_1_A_1:ForceSI}
  \lean{IPhO2026Problems.IPhO2026_1_A_1.ForceSI}
  A force read in newtons.
\end{definition}
\begin{proof}
  This is the defining declaration; unfolding it gives the stated typed data or relation.
\end{proof}

\begin{definition}[TorqueSI]
  \label{decl:physics:IPhO_2026_1_A_1:TorqueSI}
  \lean{IPhO2026Problems.IPhO2026_1_A_1.TorqueSI}
  A torque read in newton metres.
\end{definition}
\begin{proof}
  This is the defining declaration; unfolding it gives the stated typed data or relation.
\end{proof}

\begin{definition}[FigurePointLabel]
  \label{decl:physics:IPhO_2026_1_A_1:FigurePointLabel}
  \lean{IPhO2026Problems.IPhO2026_1_A_1.FigurePointLabel}
  Point labels appearing in Figure 1a.
\end{definition}
\begin{proof}
  This is the defining declaration; unfolding it gives the stated typed data or relation.
\end{proof}

\begin{definition}[WallOrientation]
  \label{decl:physics:IPhO_2026_1_A_1:WallOrientation}
  \lean{IPhO2026Problems.IPhO2026_1_A_1.WallOrientation}
  Orientations relevant to the wall in Figure 1a.
\end{definition}
\begin{proof}
  This is the defining declaration; unfolding it gives the stated typed data or relation.
\end{proof}

\begin{definition}[SubmersionStatus]
  \label{decl:physics:IPhO_2026_1_A_1:SubmersionStatus}
  \lean{IPhO2026Problems.IPhO2026_1_A_1.SubmersionStatus}
  The possible submersion status of the sealing block.
\end{definition}
\begin{proof}
  This is the defining declaration; unfolding it gives the stated typed data or relation.
\end{proof}

\begin{definition}[HingeFriction]
  \label{decl:physics:IPhO_2026_1_A_1:HingeFriction}
  \lean{IPhO2026Problems.IPhO2026_1_A_1.HingeFriction}
  The hinge property relevant to the onset of rotation.
\end{definition}
\begin{proof}
  This is the defining declaration; unfolding it gives the stated typed data or relation.
\end{proof}

\begin{definition}[AxisOrientation]
  \label{decl:physics:IPhO_2026_1_A_1:AxisOrientation}
  \lean{IPhO2026Problems.IPhO2026_1_A_1.AxisOrientation}
  The orientation of the hinge axis relative to the plane of Figure 1a.
\end{definition}
\begin{proof}
  This is the defining declaration; unfolding it gives the stated typed data or relation.
\end{proof}

\begin{definition}[ReservoirFluid]
  \label{decl:physics:IPhO_2026_1_A_1:ReservoirFluid}
  \lean{IPhO2026Problems.IPhO2026_1_A_1.ReservoirFluid}
  The material in the two reservoirs.
\end{definition}
\begin{proof}
  This is the defining declaration; unfolding it gives the stated typed data or relation.
\end{proof}

\begin{definition}[GateConfiguration]
  \label{decl:physics:IPhO_2026_1_A_1:GateConfiguration}
  \lean{IPhO2026Problems.IPhO2026_1_A_1.GateConfiguration}
  \uses{decl:physics:IPhO_2026_1_A_1:FigurePointLabel, decl:physics:IPhO_2026_1_A_1:WallOrientation, decl:physics:IPhO_2026_1_A_1:SubmersionStatus, decl:physics:IPhO_2026_1_A_1:HingeFriction, decl:physics:IPhO_2026_1_A_1:AxisOrientation, decl:physics:IPhO_2026_1_A_1:ReservoirFluid}
  Qualitative apparatus data stated in the problem and shown in Figure 1a.
\end{definition}
\begin{proof}
  This is the defining declaration; unfolding it gives the stated typed data or relation.
\end{proof}

\begin{definition}[Figure1aGeometry]
  \label{decl:physics:IPhO_2026_1_A_1:Figure1aGeometry}
  \lean{IPhO2026Problems.IPhO2026_1_A_1.Figure1aGeometry}
  \uses{decl:physics:IPhO_2026_1_A_1:LengthSI, decl:physics:IPhO_2026_1_A_1:AreaSI, decl:physics:IPhO_2026_1_A_1:VolumeSI}
  The dimensional geometric readouts from Figure 1a.
\end{definition}
\begin{proof}
  This is the defining declaration; unfolding it gives the stated typed data or relation.
\end{proof}

\begin{definition}[HydrostaticGateState]
  \label{decl:physics:IPhO_2026_1_A_1:HydrostaticGateState}
  \lean{IPhO2026Problems.IPhO2026_1_A_1.HydrostaticGateState}
  \uses{decl:physics:IPhO_2026_1_A_1:LengthSI, decl:physics:IPhO_2026_1_A_1:MassDensitySI, decl:physics:IPhO_2026_1_A_1:AccelerationSI, decl:physics:IPhO_2026_1_A_1:PressureSI, decl:physics:IPhO_2026_1_A_1:ForceSI, decl:physics:IPhO_2026_1_A_1:TorqueSI}
  Dimensional state variables used in the limiting moment balance.
\end{definition}
\begin{proof}
  This is the defining declaration; unfolding it gives the stated typed data or relation.
\end{proof}

\begin{definition}[MatchesProblemSetup]
  \label{decl:physics:IPhO_2026_1_A_1:MatchesProblemSetup}
  \lean{IPhO2026Problems.IPhO2026_1_A_1.MatchesProblemSetup}
  \uses{decl:physics:IPhO_2026_1_A_1:GateConfiguration, decl:physics:IPhO_2026_1_A_1:Figure1aGeometry, decl:physics:IPhO_2026_1_A_1:HydrostaticGateState}
  The qualitative conditions and numerical data supplied by the problem.
\end{definition}
\begin{proof}
  This is the defining declaration; unfolding it gives the stated typed data or relation.
\end{proof}

\begin{definition}[MatchesFigure1a]
  \label{decl:physics:IPhO_2026_1_A_1:MatchesFigure1a}
  \lean{IPhO2026Problems.IPhO2026_1_A_1.MatchesFigure1a}
  \uses{decl:physics:IPhO_2026_1_A_1:Figure1aGeometry}
  Metric relations read from Figure 1a.
\end{definition}
\begin{proof}
  This is the defining declaration; unfolding it gives the stated typed data or relation.
\end{proof}

\begin{definition}[HydrostaticGateLaws]
  \label{decl:physics:IPhO_2026_1_A_1:HydrostaticGateLaws}
  \lean{IPhO2026Problems.IPhO2026_1_A_1.HydrostaticGateLaws}
  \uses{decl:physics:IPhO_2026_1_A_1:Figure1aGeometry, decl:physics:IPhO_2026_1_A_1:HydrostaticGateState}
  The governing physical laws at the maximum permissible level difference: hydrostatic pressure, resultant force, buoyancy-reduced weight, the two moments about the frictionless hinge “O”, and equality of those moments at incipient rotation.
\end{definition}
\begin{proof}
  This is the defining declaration; unfolding it gives the stated typed data or relation.
\end{proof}

\begin{definition}[RoundsToNearestCentimeterSI]
  \label{decl:physics:IPhO_2026_1_A_1:RoundsToNearestCentimeterSI}
  \lean{IPhO2026Problems.IPhO2026_1_A_1.RoundsToNearestCentimeterSI}
  \uses{decl:physics:IPhO_2026_1_A_1:LengthSI}
  “length” rounds to “centimetreCount” centimetres when its SI readout lies within half a centimetre of that decimal value.
\end{definition}
\begin{proof}
  This is the defining declaration; unfolding it gives the stated typed data or relation.
\end{proof}

\begin{theorem}[Physics formalization target]
\label{thm:physics:IPhO_2026_1_A_1:target}
\lean{IPhO2026Problems.IPhO2026_1_A_1.sideLength_at_maximumLevelDifference}
\uses{decl:physics:IPhO_2026_1_A_1:LengthSI, decl:physics:IPhO_2026_1_A_1:AreaSI, decl:physics:IPhO_2026_1_A_1:VolumeSI, decl:physics:IPhO_2026_1_A_1:MassDensitySI, decl:physics:IPhO_2026_1_A_1:AccelerationSI, decl:physics:IPhO_2026_1_A_1:PressureSI, decl:physics:IPhO_2026_1_A_1:ForceSI, decl:physics:IPhO_2026_1_A_1:TorqueSI, decl:physics:IPhO_2026_1_A_1:FigurePointLabel, decl:physics:IPhO_2026_1_A_1:WallOrientation, decl:physics:IPhO_2026_1_A_1:SubmersionStatus, decl:physics:IPhO_2026_1_A_1:HingeFriction, decl:physics:IPhO_2026_1_A_1:AxisOrientation, decl:physics:IPhO_2026_1_A_1:ReservoirFluid, decl:physics:IPhO_2026_1_A_1:GateConfiguration, decl:physics:IPhO_2026_1_A_1:Figure1aGeometry, decl:physics:IPhO_2026_1_A_1:HydrostaticGateState, decl:physics:IPhO_2026_1_A_1:MatchesProblemSetup, decl:physics:IPhO_2026_1_A_1:MatchesFigure1a, decl:physics:IPhO_2026_1_A_1:HydrostaticGateLaws, decl:physics:IPhO_2026_1_A_1:RoundsToNearestCentimeterSI}
The assigned autoformalize agent should translate this physics problem into Lean declarations in the covered file, with theorem and lemma proof bodies written as `by sorry`.
\end{theorem}
\begin{proof}
This is an autoformalization task, not a proof task. Produce statements that can later be proved by the physics prover without weakening the physical model.
\end{proof}
% --- Archon physics formalization source end ---
